\section{引言}

张量分解为多维数据提供了自然的低秩表示方式，已广泛应用于图像恢复、信号处理、压缩表示与缺失补全等任务\cite{Kolda2009Tensor,DeLathauwer2000MLSVD,Cichocki2015SignalProcessing,Sidiropoulos2017TSP,Liu2013TensorCompletion}。作为经典模型之一，Tucker 分解通过核心张量与模因子矩阵刻画多模之间的耦合关系，在结构可解释性、秩控制灵活性与算法可实现性之间取得了良好平衡\cite{tucker1966some, de2000multilinear, kolda2009tensor}。然而，标准 Tucker 分解本质上建立在多线性条目生成假设之上；当观测张量来源于非线性生成、饱和响应、复杂高阶交互或异方差扰动时，纯多线性模型往往会产生系统性失配，进而限制其在重建、恢复与后续预测任务中的表现。

为缓解这一问题，近年来出现了多种非线性张量分解方法。现有工作一方面通过链接函数、广义似然、核方法或高斯过程改变潜在线性预测到观测条目的映射，另一方面也通过深度先验、神经解码器或网络化分解结构提升模型表达力 \cite{hong2020generalized, xu2011infinite, saragadam2024deeptensor, varolgunes2023nmtucker, zhai2024quadratic}。 这些方法表明，非线性扩展确实能够显著增强张量模型对复杂数据机制的适应性；但与此同时，它们通常也伴随更高的推断与训练复杂度、更大的超参数设计空间，或对经典 Tucker 结构透明性的削弱。因而，一个仍然值得回答的问题是：能否在尽量保留 Tucker 低秩骨架、参数效率与多模可解释性的前提下，为其引入受控而有效的非线性表达能力？ 这也是本文的出发点。

本文提出多项式链接的非线性 Tucker 分解（Nonlinear Tucker Decomposition with Polynomial Links, NTD-PL）。该模型以共享的 Tucker 潜在张量
\[
\cS = \cX \times_1 \mathbf{A}^{(1)} \times_2 \cdots \times_N \mathbf{A}^{(N)}
\]
作为低秩 骨架，并在其上施加显式的逐元素多项式链接
\[
\hat{\cY} = f_{\beta}(\cS) = \sum_{p=0}^{P} \beta_p \cS^{\odot p}.
\]
由此，标准 Tucker 被自然包含为线性特例，而高阶非线性则通过少量链接系数以受控方式引入。与依赖核/高斯过程的非参数映射不同，NTD-PL 的非线性形式是显式、低维且确定性的；与依赖神经解码器的深度非线性模型不同，NTD-PL 并不替换 Tucker 的低秩骨架，而是仅对条目生成机制进行结构化增强。因而，它更适合被理解为一种增强 Tucker 的模型，而非对 Tucker 的完全替代。

这一建模方式的关键在于，它并不是简单地对 Tucker 重构结果施加逐元素后处理。由于潜在张量$\cS$的每个条目本身已经由核心张量与各模因子的多线性耦合生成，因此$\cS^{\odot p}$对应的是共享低秩骨架上的高阶交互重组，而不是独立于低秩结构之外的外部校正。正因如此，NTD-PL 在不显式引入高阶因子张量的情况下，仍能够获得超出纯多线性模型的表达能力，同时保持与 Tucker 同量级的参数化骨架。

围绕该模型，本文进一步刻画了其若干基本性质。首先，NTD-PL 与标准 Tucker 构成清晰的嵌套关系：当高阶系数关闭时，模型精确退化为 Tucker；随着多项式次数增加，模型类按次数形成受控扩张。其次，在观测张量由逐元素多项式函数作用于某个 Tucker 潜在张量生成时，NTD-PL 可在相应阶数下实现精确匹配；当生成函数为解析函数时，有限阶多项式链接可给出系统逼近，并具有显式截断误差界。再次，为了同时适用于完整张量分解与带掩码观测的张量补全，本文构建了统一的学习目标，并结合系数 ridge 更新、因子归一化与次数延续策略给出可实现的优化框架。

实验结果进一步支持了上述建模与理论分析。在可控合成实验中，NTD-PL 在严格线性数据上与 Tucker 保持一致，并在多项式与平滑非线性残差下随非线性强度增大持续优于 Tucker、CP、TT 和 TR 等基线，且其性能提升与目标函数的主导阶次相一致。在真实高光谱数据上，NTD-PL 在 CAVE 数据集的全观测重构与随机缺失补全任务中均稳定优于 Tucker，这种增益不仅体现在场景均值上，也在多数场景、多个缺失率和局部空间区域上保持一致。进一步的机制分析表明，该优势并不能由对 Tucker 输出施加简单多项式后处理所替代，而主要来源于显式非线性链接与共享低秩骨架的联合建模；跨数据集实验则说明，该方法在多数真实场景中有效，同时也呈现出清晰的适用边界。本文的主要贡献概括如下。

\begin{enumerate}
    \item 提出 NTD-PL，一种在共享 Tucker 潜在张量上施加显式逐元素多项式链接的非线性低秩张量模型。
    \item 说明该模型与 Tucker 的嵌套关系、参数规模特征，以及其在多项式生成和解析函数生成情形下的表达性质。
    \item 给出统一的带掩码学习形式，使模型能够同时处理完整张量分解与张量补全。
    \item 通过可控合成实验与真实高光谱实验验证了该模型的有效性与适用边界。
\end{enumerate}

本文其余部分安排如下：第二节综述相关工作；第三节给出记号与预备知识；第四节介绍模型、性质与优化方法；第五节报告实验结果；第六节总结全文并讨论后续工作。
